\chapter{實驗設計}
\label{ch:experiments}

\section{Benchmark and metric}

本文主要使用 LIBERO-Spatial 作為 simulation benchmark~\cite{liu2023libero}。Spatial tasks 對空間關係、gripper-object alignment 與局部接觸幾何較敏感，因此適合觀察 action chunk stale execution 的影響。若 frozen planner 每一步都重新觀察並重新規劃，local residual 不一定有幫助；但當 chunk 後段 action 在較舊的 visual feedback 下被執行時，wrist-based correction 才有可能展現價值。

主要 metric 是 success rate：

\begin{equation}
  \mathrm{SuccessRate}=\frac{\mathrm{Number\ of\ successful\ episodes}}{\mathrm{Number\ of\ evaluated\ episodes}}\times100\%.
\end{equation}

本文以 success rate 作為主結果。Figures 中若出現 Wilson interval，僅用於視覺化 episode-level uncertainty；本文不把它解讀為 multi-seed confidence interval，也不以 margin plots 取代 success-rate reporting。

\section{Registry-backed evaluation}

所有 quantitative claims 都由 project eval registry 追蹤。Registry 的角色不是提供新的 metric，而是確保每個 success rate 都能追溯到對應的 policy、benchmark split、chunk protocol、episode count 與 status。這避免了兩種常見錯誤：第一，把不同 chunk protocol 的結果直接比較；第二，從 terminal log 手動抄數字而失去 provenance。Exact registry file paths 與 source manifest 列於 Appendix~\ref{app:thesis-implementation-reproducibility}。

\section{Protocol variables}

本文所有主要比較都控制三個 chunk protocol variables：

\begin{table}[htbp]
  \centering
  \caption{Chunk protocol variables。}
  \label{tab:protocol_vars}
  \begin{tabular}{ll}
    \toprule
    Protocol variable & Meaning \\
    \midrule
    Planned chunk length & Slow planner 一次產生多少 candidate actions \\
    Executed open-loop length & Robot 在下一次切換前連續執行多少 actions \\
    Replanning interval & Slow planner 多久接收新 observation 並 refresh chunk \\
    \bottomrule
  \end{tabular}
\end{table}

若 \method{} 與 \smolvla{} baseline 使用不同 protocol，success rate 不能視為 matched comparison。因此本文的 main sweeps 固定兩者使用相同 chunk override，只改變 protocol 本身或 residual calibration setting。

\section{Synchronous sweeps}

除了 planner-delay stress test 以外，本文所有主要 sweeps 都使用 synchronous inference。這代表 slow planner 在需要新 chunk 時立即產生可用 chunk；實驗重點不是 planner compute latency，而是 action chunk 被 open-loop 執行時產生的 stale feedback problem。

\subsection{Plan-50 execution/replan sweep}

第一個 sweep 固定 planned chunk length 為 50，並令 executed open-loop length 與 replanning interval 共同等於 $K$。此設定測試：當 slow planner 產生長 chunk 時，若 robot 每 $K$ 步才切換或 refresh chunk，fast wrist residual 是否能改善後段 stale actions。

這個 protocol 特別適合回答「長 planning horizon 下，是否可以用 fast local correction 延緩 stale chunk degradation」。它不代表 full async scheduler，因為 planner 在每次 request 時仍被視為 synchronous available。

\subsection{Matched chunk sweep}

第二個 sweep 將 planned chunk length、executed open-loop length 與 replanning interval 全部設為同一個 $K$。此設定比 plan-50 sweep 更嚴格，因為 planner 不再固定產生 50-step chunk；每個 tested $K$ 都是一個 matched planning/execution regime。

Matched chunk sweep 的目的，是確認 \method{} 的改善不是來自「長 chunk 但頻繁 refresh」這種混合設定，而是在真正 matched long-chunk execution 中仍能降低 degradation。

\section{Residual calibration}

Residual calibration sweep 改變 $\alpha$ 與 $\dmax$，觀察 residual authority 對 success rate 的影響。這組實驗回答的問題是：fast residual 應該是有界 local correction，還是可以讓 residual head 自由輸出大幅 action correction。

本文預期的合理區間是 bounded correction。若 $\alpha$ 或 $\dmax$ 過大，residual 可能覆蓋 frozen planner 的 global intent；若過小，則無法修正 chunk 後段累積的局部偏差。因此 calibration sweep 不只是 hyperparameter search，而是驗證 clipped merge 是否為方法必要成分。

\section{Async-timestep planner-delay stress test}

Planner-delay sweep 是本文唯一 async-timestep inference experiment。此實驗中，slow planner 在 step $t$ 發出 request，但 chunk 在 $t+d$ 個 control timesteps 後才可用；execution 從 delayed chunk 的對應 index 開始。Fast wrist residual 仍在 current timestep 執行，因此能利用當前 wrist observation 修正 delayed base action。

這個實驗和 synchronous sweeps 分開報告，因為它測試的是不同 runtime assumption。它不是完整 real-time remote-planner deployment，也不涵蓋任意 compute delays；它是一個受控 stress test，用來檢查 slow chunks late arrival 時，current wrist evidence 是否能降低 disable-fast ablation 的 degradation。

\section{Unfinished experiments}

以下項目目前不作為 main-text evidence：

\begin{itemize}
  \item Dual-view DINO fast residual ablation：同時使用 wrist 與 third-person image 的 fast model。
  \item Latent-context ablation：移除 cached slow-planner context。
  \item Chunk-age/stale-target training：用 delayed base chunks 訓練 residual targets。
  \item No-residual-clip control sweeps：擴展 no-clip controls。
  \item Multi-seed confidence intervals：以多 seeds 或 task-order variants 重跑 main sweeps。
  \item Causal visual ablations：wrist mask、third-person mask、DINO patch masking。
  \item Direct A2C2 same-backbone comparison：在相同 \smolvla{} backbone 與 registry 下重現 A2C2。
  \item Real-robot tabletop validation：在 physical tabletop setup 執行 matched trials。
\end{itemize}
